\documentclass[11pt,a4paper]{article}

\usepackage[utf8]{inputenc}
\usepackage[T1]{fontenc}

\usepackage[margin=2cm]{geometry}
\usepackage{booktabs}
\usepackage{listings}
\usepackage{xcolor}
\usepackage{titlesec}
\usepackage{enumitem}

\lstset{
    basicstyle=\ttfamily\small,
    breaklines=true,
    frame=single,
    backgroundcolor=\color{gray!10},
    keywordstyle=\color{blue},
    commentstyle=\color{gray},
    stringstyle=\color{red!70!black},
    showstringspaces=false
}

\titleformat{\section}{\large\bfseries\color{blue!70!black}}{}{0em}{}[\titlerule]
\titleformat{\subsection}{\normalsize\bfseries}{}{0em}{}

\title{\textbf{Swift Safety Patterns}\\[0.3em]
\large Eliminate Force-Unwraps \& Force-Casts in Production Code}
\author{Senior Engineering Review — May 2026}
\date{}

\begin{document}

\maketitle

\begin{abstract}
\noindent\textbf{Goal:} Remove all unsafe \texttt{!} (force-unwrap) and \texttt{as!} (force-cast) patterns.\\
\textbf{Philosophy:} Prefer safety and clarity over brevity.
\end{abstract}

\section{1. Property Access on Optionals}

\begin{tabular}{@{}p{5.5cm}p{5.5cm}l@{}}
\toprule
\textbf{Bad (Avoid)} & \textbf{Good (Preferred)} & \textbf{Notes} \\
\midrule
\texttt{foo!.bar = 42} & \texttt{foo?.bar = 42} & Default choice \\
\texttt{constraint.constant = 100} & \texttt{constraint?.constant = 100} & Auto Layout constraints \\
\texttt{pendingRefresh!.cancel()} & \texttt{pendingRefresh?.cancel()} & Optional work items \\
\bottomrule
\end{tabular}

\section{2. Creating URLs from Strings}

\textbf{Bad:}
\begin{lstlisting}
let url = URL(string: "https://example.com")!
\end{lstlisting}

\textbf{Good:}
\begin{lstlisting}
let url = makeURL("https://example.com")
\end{lstlisting}

\textbf{Recommended Helper:}
\begin{lstlisting}[language=Swift]
/// Creates a URL or crashes with a clear message during development.
private static func makeURL(_ string: String) -> URL {
    guard let url = URL(string: string) else {
        fatalError("Invalid hardcoded URL: \(string)")
    }
    return url
}
\end{lstlisting}

\section{3. Calendar \& Date Calculations}

\textbf{Bad:}
\begin{lstlisting}
let nextDate = Calendar.current.date(byAdding: .minute, value: 15, to: Date())!
\end{lstlisting}

\textbf{Good:}
\begin{lstlisting}
let nextDate = Calendar.current.date(byAdding: .minute, value: 15, to: Date())
    ?? Date().addingTimeInterval(15 * 60)
\end{lstlisting}

\section{4. Force Casts (\texttt{as!})}

\textbf{Bad:}
\begin{lstlisting}
let cell = collectionView.dequeueReusableCell(...) as! LanguageCell
\end{lstlisting}

\textbf{Good:}
\begin{lstlisting}
guard let cell = collectionView.dequeueReusableCell(
    withReuseIdentifier: "LanguageCell", 
    for: indexPath
) as? LanguageCell else {
    fatalError("Failed to dequeue LanguageCell")
}
\end{lstlisting}

\section{5. Implicitly Unwrapped Optionals (Declaration)}

\textbf{Bad:}
\begin{lstlisting}
private var speakerImageHeightConstraint: NSLayoutConstraint!
\end{lstlisting}

\textbf{Good:}
\begin{lstlisting}
private var speakerImageHeightConstraint: NSLayoutConstraint?
\end{lstlisting}

\section{6. DispatchWorkItem \& Closures}

\textbf{Bad:}
\begin{lstlisting}
DispatchQueue.main.asyncAfter(..., execute: pendingRefresh!)
\end{lstlisting}

\textbf{Good:}
\begin{lstlisting}
let workItem = DispatchWorkItem { ... }
pendingRefresh = workItem
DispatchQueue.main.asyncAfter(..., execute: workItem)
\end{lstlisting}

\section{7. Safe Array Access}

\textbf{Bad:}
\begin{lstlisting}
let item = array[0]
\end{lstlisting}

\textbf{Good:}
\begin{lstlisting}
let item = array.first ?? array[0]   // final fallback only
\end{lstlisting}

\section{8. Timestamp / Throttle Checks}

\textbf{Bad:}
\begin{lstlisting}
Date().timeIntervalSince(lastServerSelectionTime!)
\end{lstlisting}

\textbf{Good:}
\begin{lstlisting}
if let last = lastServerSelectionTime, 
   Date().timeIntervalSince(last) > 10.0 {
\end{lstlisting}

\section{When \texttt{!} Is Still Acceptable}

Use force-unwrap \textbf{only} in these cases:

\begin{itemize}[leftmargin=*]
    \item Immediately after a successful \texttt{guard let} or \texttt{if let}
    \item In unit tests (where crashing on failure is often desired)
    \item With a clear explanatory comment
\end{itemize}

\textbf{Example:}
\begin{lstlisting}
guard let constraint = speakerImageHeightConstraint else { return }
constraint.constant = 100     // Safe here — no need for !
\end{lstlisting}

\section{Final Rule of Thumb}

\begin{center}
\fbox{\parbox{0.9\textwidth}{\centering
\textbf{Default to \texttt{?}} (optional chaining)\\[0.5em]
Only use \texttt{!} when you have \textbf{already proven} the value exists.
}}
\end{center}

\vspace{1em}
\hrule
\vspace{0.5em}
\footnotesize
\textit{Maintained by Senior Engineering — Lutheran Radio Project, May 2026}

\end{document}
